% GENERATED FILE. Edit canonical lessons, then run npm run generate:books.
% canonical-source-hash: fnv1a64:78ea44e4a352c453
% canonical-lessons: JA-C343-shin-you-suru, JA-C343-shinrai, JA-C343-kabau, JA-C343-bareru, JA-C343-kimetsukeru

\chapter{Trust and Doubt}
\label{ch:ja-trust-and-doubt}
\hlchaptermodality{343}

\begin{chapteropening}
\textbf{By the end of this chapter:} \emph{I can say whom I trust and why.}

Complete the last lesson of chapter 343: I can say whom I trust and why.
\end{chapteropening}

\section[\texorpdfstring{\ja{しんようする} (shin'yō suru)}{しんようする (shin'yō suru)}]{\ja{しんようする} (shin'yō suru) — to trust; to believe in}

\label{lesson:JA-C343-shin-you-suru}



\begin{quote}
Before the new one: say the Japanese for the reason is, then the Japanese for an attitude.
\end{quote}

\subsection*{You'll want to know: \ja{しんようする}}
\textbf{\ja{しんようする}} — \emph{shin'yō suru} — \textquotedblleft{}to trust; to believe in\textquotedblright{}.

\begin{quote}
\textbf{\ja{うそを} \ja{言}\ja{ったので、かれを} \ja{しんようしません。}} — \emph{uso o itta node, kare o shin'yō shimasen} — \textquotedblleft{}He lied, so I don't trust him.\textquotedblright{}
\end{quote}

Polite \textbf{\ja{しんようします}}, \emph{shin'yōshimasu}; a \ja{する} verb. Past (polite): \textbf{\ja{みんなが} \ja{かれを} \ja{しんようしました。}}, \emph{minna ga kare o shin'yō shimashita}.

\begin{cousinweb}
\textbf{\ja{しんようする}} (\emph{shin'yō suru}) is built from a Sino-Japanese noun and the native verb \emph{suru}, \textquotedblleft{}to do\textquotedblright{}. The noun, \emph{shin'yō}, is two Chinese characters read with their on'yomi, readings borrowed from Chinese: \emph{shin}, \textquotedblleft{}trust\textquotedblright{} (Middle Chinese \emph{sinH}) and \emph{yō}, \textquotedblleft{}utilize\textquotedblright{} (\emph{yowngH}) (KANJIDIC2; JMdict; Middle Chinese in Baxter's transcription).
\end{cousinweb}

\begin{grammarlens}[title={What you've built}]
One more word for giving reasons and opinions.
\end{grammarlens}

\subsection*{Your turn}
\begin{itemize}
\raggedright
  \item \emph{Say it:} \emph{shin'yō suru}
  \item \emph{Say it:} \emph{shin'yō suru}, once more
  \item \emph{Say it:} \emph{uso o itta node, kare o shin'yō shimasen} — He lied, so I don't trust him.
  \item \emph{Recall:} say the Japanese for the reason is, then the Japanese for an attitude, then say \emph{shin'yō suru} again
\end{itemize}

\begin{tcolorbox}[breakable,colback=teal!4,colframe=teal!35!black,title={Before you move on}]
What does \ja{しんようする} mean? (\textquotedblleft{}To trust; to believe in\textquotedblright{}.) Say it once more.
\end{tcolorbox}
\section[\texorpdfstring{\ja{しんらい} (shinrai)}{しんらい (shinrai)}]{\ja{しんらい} (shinrai) — trust; confidence}

\label{lesson:JA-C343-shinrai}



\begin{quote}
Before the new one: say the Japanese for an attitude, then the Japanese for to trust.
\end{quote}

\subsection*{You'll want to know: \ja{しんらい}}
\textbf{\ja{しんらい}} — \emph{shinrai} — \textquotedblleft{}trust; confidence\textquotedblright{}.

\begin{quote}
\textbf{\ja{うそを} \ja{言}\ja{ったら、しんらいを} \ja{なくします。}} — \emph{uso o ittara, shinrai o nakushimasu} — \textquotedblleft{}If you lie, you lose people's trust.\textquotedblright{}
\end{quote}

\begin{cousinweb}
\textbf{\ja{しんらい}} (\emph{shinrai}) is Sino-Japanese (kango): two Chinese characters read with their on'yomi, readings borrowed from Chinese. They are \emph{shin}, \textquotedblleft{}trust\textquotedblright{} (Middle Chinese \emph{sinH}) and \emph{rai}, \textquotedblleft{}trust\textquotedblright{} (\emph{lajH}) (KANJIDIC2; Middle Chinese in Baxter's transcription).
\end{cousinweb}

\begin{grammarlens}[title={What you've built}]
One more word for giving reasons and opinions.
\end{grammarlens}

\subsection*{Your turn}
\begin{itemize}
\raggedright
  \item \emph{Say it:} \emph{shinrai}
  \item \emph{Say it:} \emph{shinrai}, once more
  \item \emph{Say it:} \emph{uso o ittara, shinrai o nakushimasu} — If you lie, you lose people's trust.
  \item \emph{Recall:} say the Japanese for an attitude, then the Japanese for to trust, then say \emph{shinrai} again
\end{itemize}

\begin{tcolorbox}[breakable,colback=teal!4,colframe=teal!35!black,title={Before you move on}]
What does \ja{しんらい} mean? (\textquotedblleft{}Trust; confidence\textquotedblright{}.) Say it once more.
\end{tcolorbox}
\section[\texorpdfstring{\ja{かばう} (kabau)}{かばう (kabau)}]{\ja{かばう} (kabau) — to stand up for (someone)}

\label{lesson:JA-C343-kabau}



\begin{quote}
Before the new one: say the Japanese for to trust, then the Japanese for confidence.
\end{quote}

\subsection*{You'll want to know: \ja{かばう}}
\textbf{\ja{かばう}} — \emph{kabau} — \textquotedblleft{}to stand up for (someone)\textquotedblright{}.

\begin{quote}
\textbf{\ja{ちちが} \ja{おこったので、あには} \ja{おとうとを} \ja{かばいました。}} — \emph{chichi ga okotta node, ani wa otōto o kabaimashita} — \textquotedblleft{}My father got angry, so my older brother stood up for my younger brother.\textquotedblright{}
\end{quote}

Polite \textbf{\ja{かばいます}}, \emph{kabaimasu}; a godan verb. Past (polite): \textbf{\ja{あねが} \ja{いもうとを} \ja{かばいました。}}, \emph{ane ga imōto o kabaimashita}.

\begin{cousinweb}
\textbf{\ja{かばう}} (\emph{kabau}) is a native Japanese word (wago), inherited rather than borrowed from Chinese, and its older origin is not settled. JMdict notes that it is usually written in kana (JMdict).
\end{cousinweb}

\begin{grammarlens}[title={What you've built}]
One more word for giving reasons and opinions.
\end{grammarlens}

\subsection*{Your turn}
\begin{itemize}
\raggedright
  \item \emph{Say it:} \emph{kabau}
  \item \emph{Say it:} \emph{kabau}, once more
  \item \emph{Say it:} \emph{chichi ga okotta node, ani wa otōto o kabaimashita} — My father got angry, so my older brother stood up for my younger brother.
  \item \emph{Recall:} say the Japanese for to trust, then the Japanese for confidence, then say \emph{kabau} again
\end{itemize}

\begin{tcolorbox}[breakable,colback=teal!4,colframe=teal!35!black,title={Before you move on}]
What does \ja{かばう} mean? (\textquotedblleft{}To stand up for (someone)\textquotedblright{}.) Say it once more.
\end{tcolorbox}
\section[\texorpdfstring{\ja{ばれる} (bareru)}{ばれる (bareru)}]{\ja{ばれる} (bareru) — to come out (of a lie or secret); to be found out}

\label{lesson:JA-C343-bareru}



\begin{quote}
Before the new one: say the Japanese for confidence, then the Japanese for to stand up for.
\end{quote}

\subsection*{You'll want to know: \ja{ばれる}}
\textbf{\ja{ばれる}} — \emph{bareru} — \textquotedblleft{}to come out (of a lie or secret); to be found out\textquotedblright{}.

\begin{quote}
\textbf{\ja{うそが} \ja{ばれたので、ちちは} \ja{おこりました。}} — \emph{uso ga bareta node, chichi wa okorimashita} — \textquotedblleft{}The lie came out, so my father got angry.\textquotedblright{}
\end{quote}

Polite \textbf{\ja{ばれます}}, \emph{baremasu}; an ichidan verb. Past (polite): \textbf{\ja{うそが} \ja{ばれました。}}, \emph{uso ga baremashita}.

\begin{cousinweb}
\textbf{\ja{ばれる}} (\emph{bareru}) is a native Japanese word (wago), inherited rather than borrowed from Chinese, and its older origin is not settled; JMdict gives it no kanji spelling (JMdict).
\end{cousinweb}

\begin{grammarlens}[title={What you've built}]
One more word for giving reasons and opinions.
\end{grammarlens}

\subsection*{Your turn}
\begin{itemize}
\raggedright
  \item \emph{Say it:} \emph{bareru}
  \item \emph{Say it:} \emph{bareru}, once more
  \item \emph{Say it:} \emph{uso ga bareta node, chichi wa okorimashita} — The lie came out, so my father got angry.
  \item \emph{Recall:} say the Japanese for confidence, then the Japanese for to stand up for, then say \emph{bareru} again
\end{itemize}

\begin{tcolorbox}[breakable,colback=teal!4,colframe=teal!35!black,title={Before you move on}]
What does \ja{ばれる} mean? (\textquotedblleft{}To come out (of a lie or secret); to be found out\textquotedblright{}.) Say it once more.
\end{tcolorbox}
\section[\texorpdfstring{\ja{きめつける} (kimetsukeru)}{きめつける (kimetsukeru)}]{\ja{きめつける} (kimetsukeru) — to jump to conclusions; to decide one-sidedly}

\label{lesson:JA-C343-kimetsukeru}



\begin{quote}
Before the new one: say the Japanese for to stand up for, then the Japanese for to be found out.
\end{quote}

\subsection*{You'll want to know: \ja{きめつける}}
\textbf{\ja{きめつける}} — \emph{kimetsukeru} — \textquotedblleft{}to jump to conclusions; to decide one-sidedly\textquotedblright{}.

\begin{quote}
\textbf{\ja{まだ} \ja{わからないから、きめつけないで} \ja{ください。}} — \emph{mada wakaranai kara, kimetsukenaide kudasai} — \textquotedblleft{}We don't know yet, so please do not jump to conclusions.\textquotedblright{}
\end{quote}

Polite \textbf{\ja{きめつけます}}, \emph{kimetsukemasu}; an ichidan verb. Past (polite): \textbf{\ja{ちちが} \ja{きめつけました。}}, \emph{chichi ga kimetsukemashita}.

\begin{cousinweb}
\textbf{\ja{きめつける}} (\emph{kimetsukeru}) is a native Japanese word (wago), inherited rather than borrowed from Chinese: it is built from \emph{kime}, \textquotedblleft{}deciding (kimeru)\textquotedblright{} and \emph{tsukeru}, \textquotedblleft{}to attach\textquotedblright{}, and its kanji were attached later as kun'yomi (JMdict).
\end{cousinweb}

\begin{grammarlens}[title={What you've built}]
One more word for giving reasons and opinions.
\end{grammarlens}

\subsection*{Your turn}
\begin{itemize}
\raggedright
  \item \emph{Say it:} \emph{kimetsukeru}
  \item \emph{Say it:} \emph{kimetsukeru}, once more
  \item \emph{Say it:} \emph{mada wakaranai kara, kimetsukenaide kudasai} — We don't know yet, so please do not jump to conclusions.
  \item \emph{Recall:} say the Japanese for to stand up for, then the Japanese for to be found out, then say \emph{kimetsukeru} again
\end{itemize}

\begin{tcolorbox}[breakable,colback=teal!4,colframe=teal!35!black,title={Before you move on}]
What does \ja{きめつける} mean? (\textquotedblleft{}To jump to conclusions; to decide one-sidedly\textquotedblright{}.) Say it once more.
\end{tcolorbox}
